\bookchapter{פרק 16: מערכת חשבונות ושירותי ענן}

\chapterintro{הרבה אנשים חושבים שהחשבון משמש רק להורדת אפליקציות. למעשה, מה שחשבון טלפון אחד אחראי לו הוא הרבה יותר ממה שתדמיין:}

\booksection{אנשי קשר}

: מספרי טלפון ששמרת במשך שנים או עשרות שנים.

\begin{itemize}
  \item \textbf{אנשי קשר}: מלידת הילד ועד טקס הסיום, מהדייט הראשון ועד טיול החתונה.
  \item \textbf{פתקים}: רעיונות שרשמת במהירות, יומן, סיסמאות (כן, הרבה אנשים שומרים סיסמאות בפתקים).
  \item \textbf{פתקים}: אפליקציות בתשלום ורכישות בתוך אפליקציה שעלו מאות.
  \item \textbf{איתור הטלפון}: אחרי שהטלפון אבד, זהו הכלי היחיד שיכול לעזור לך לאתר, לנעול, למחוק נתונים מרחוק.
  \item \textbf{אם אתה משתמש באייפון, כל הדברים הללו קשורים ל-Apple ID שלך. אם אתה משתמש באנדרואיד, הם קשורים לחשבון Huawei/Xiaomi/OPPO/vivo שלך, וגם לחשבון Google שלך. החשבונות האלה הם "תעודת הזהות" של החיים הדיגיטליים שלך, אם תאבד את תעודת הזהות, הרבה דברים לא יהיו ניתנים להוכחה.}שתיים, שלושה דברים שחובה לעשות, עכשיו
\end{itemize}

הפעל אימות דו-שלבי

\booksection{: בין אם אתה משתמש ב-Apple או באנדרואיד, הדבר הראשון הוא ללכת להגדרות החשבון ולהפעיל “אימות דו-שלבי” או “אימות דו-גורמי”. זה כמו להוסיף מנעול נוסף לדלת - סיסמה לבד לא מספיקה, צריך גם קוד אימות לטלפון. גם אם מישהו ינחש את הסיסמה שלך באינטרנט, בלי הטלפון שלך הוא לא יכנס.}

\begin{enumerate}
  \item \textbf{קשר שיטת שחזור אמינה}: אימות דו-שלבי בטוח מאוד, אבל יש סיכון - מה אם הטלפון שלך עצמו אבד? לכן המערכת תציע לך להגדיר "מפתח שחזור" או "אימייל/טלפון חלופי". זהו הנתיב היחיד לשחזר את החשבון במקרה חירום, עדיף לרשום על נייר, לא לשמור בטלפון.
  \item \textbf{אשר ש"איתור הטלפון" מופעל}: פונקציה זו יכולה לעזור לך לאתר, לנעול, להשמיע צליל, ואפילו למחות מרחוק את כל הנתונים אחרי אובדן. זו קו ההגנה האחרון. הרבה אנשים מדלגים על השלב הזה בקניית הטלפון, וכשרוצים להשתמש מגלים שלא הפעילו - ואז מאוחר מדי.
  \item \textbf{שלוש, iCloud: "האוויר" של משתמשי Apple}מי שמשתמש באייפון, מוקדם או מאוחר יתמודד עם שאלה: מה עושים כשנגמר מקום ב-iCloud?
\end{enumerate}

\booksection{Apple נותנת לך 5GB בחינם, להגיד את האמת, 5GB לא מספיקים אפילו לגיבוי אחד. לכן רוב האנשים בסוף מקבלים את החלון המעצבן: “שטח האחסון ב-iCloud מלא”. בשלב זה יש לך שלוש אפשרויות:}

ניקוי

: לך להגדרות השם שלך iCloud ניהול שטח אחסון, ראה מי תופס את המקום שלך. בדרך כלל, תמונות וגיבויים הם העיקר. אפשר למחוק קודם גיבויים ישנים לא נחוצים.

\begin{itemize}
  \item \textbf{שדרוג}: 50GB ב-6 יואן לחודש, 200GB ב-21 יואן לחודש. לדעתי 50GB מספיקים לרוב האנשים, בעיקר לגיבוי ולתמונות ממוזערות.
  \item \textbf{עקיפה}: אם לא רוצה לשלם, אפשר לייצא תמונות למחשב מדי פעם, ואז לכבות סנכרון תמונות iCloud. אבל אז תאבד את הנוחות של גיבוי אוטומטי.
  \item \textbf{תזכורת חשובה: ספריית תמונות iCloud אינה "כונן רשת" - היא "סנכרון". אם תמחק תמונה מהטלפון, היא תימחק גם מ-iCloud. אם מה שאתה רוצה הוא "גיבוי לענן ואז פינוי הטלפון", אז תמונות iCloud לא הדבר שאתה צריך, אתה צריך גיבוי ענן של iCloud, או כונן רשת של צד שלישי.}ארבע, סנכרון וגיבוי - אל תתבלבל
\end{itemize}

שתי המילים נשמעות דומות, אבל ההשלכות שונות לחלוטין:

\begin{tcbinfo}{iCloud+: פרטיות בתשלום}
כל רמה בתשלום (50/200/2TB) =\textbf{iCloud+}משתמשים, פתחו שלוש פונקציות הגנת פרטיות:

\begin{itemize}
  \item \textbf{Private Relay (ממסר פרטי)}:"VPN" של Apple. Safari דרך Apple, איש (גם Apple) לא רואה הכל. ל-Wi-Fi ציבורי.\texttt{הגדרות > שמך > iCloud > Private Relay}.
  \item \textbf{Hide My Email (הסתר דוא"ל)}:כתובת אקראית חד-פעמית, העברה לדוא"ל אמיתי. סגירה = הפסקה. דוא"ל לא נחשף.\texttt{הגדרות > שמך > iCloud > Hide Email}.
  \item \textbf{דומיין דוא"ל מותאם אישית}:דומיין משלך, למותג/משפחה.
\end{itemize}

טיפ: Private Relay לא זמין בכל מקום.

\end{tcbinfo}

\begin{tcbwarn}{iCloud ADP - E2EE אמיתי}
בסוף 2022 Apple השיקה\textbf{iCloud Advanced Data Protection (ADP)}, E2EE מ-בריאות/סיסמאות ל-\textbf{רוב נתוני iCloud}—תמונות, פתקים, גיבוי, מכשיר - 23 סוגים.

\begin{itemize}
  \item \textbf{מה זה E2EE}:מפתח במכשיר, לא ב-Apple. גם אם נפרץ/חוק, Apple לא יכול לפענח.
  \item \textbf{נתיב ההפעלה}:\texttt{הגדרות > שמך > iCloud > ADP}, לפי ההנחיות הגדר איש קשר לשחזור/מפתח שחזור (\textbf{חובה}, אם תאבד את המפתח לא ניתן יהיה לשחזר את הנתונים שלך).
  \item \textbf{הערות חשובות}:1) שמור קוד שחזור; 2) אלבום משותף לא E2EE; 3) סין לא תומך.
\end{itemize}

טיפ: אם יש פרטיות חשובה והאזור תומך, הפעל ADP - ההגנה החזקה ביותר לענן.

\end{tcbwarn}

\booksection{סנכרון}

הוא דו-כיווני: אם תמחק מהטלפון, יימחק גם מהענן; אם תמחק מהענן, יימחק גם מהטלפון. תמונות iCloud, אנשי קשר, פתקים - כל אלה הם סנכרון. היתרון: עקביות בין מכשירים. החיסרון: “נמחק = נעלם לתמיד”.

\begin{itemize}
  \item \textbf{גיבוי}הוא חד-כיווני: הטלפון מצלם "תמונת מצב" של המצב הנוכחי ושומר בענן. אחר כך אם תמחק משהו מהטלפון, העותק בגיבוי הענן נשאר. גיבוי ענן iCloud, גיבוי מקומי במחשב - שניהם גיבוי. היתרון: הנתונים בטוחים יותר. החיסרון: תופס מקום, ולא מתעדכן בזמן אמת.
  \item \textbf{חמש, מטריצת החשבונות של משתמשי אנדרואיד}המצב בצד אנדרואיד מורכב יותר מ-Apple, כי הנתונים שלך עשויים להיות מפוזרים במספר חשבונות:
\end{itemize}

\booksection{חשבון מותג}

(Huawei/Xiaomi/OPPO/vivo): מנהל את שטח הענן שלך, איתור טלפון, חנות אפליקציות, ארנק. לכל מותג יש שירות ענן משלו, בדרך כלל מודל בסיסי בחינם + תשלום להרחבה.

\begin{itemize}
  \item \textbf{חשבון Google}: אם אתה משתמש ב-Gmail, Google תמונות, Google אנשי קשר, הדברים האלה נמצאים בחשבון Google. Google נותנת 15GB בחינם, יותר נדיבה מ-Apple, אבל גם משותפת עם Gmail ו-Driver.
  \item \textbf{WeChat/QQ}: הם בלתי תלויים במערכת הטלפון. היסטוריית הצ'אט והמועדפים של WeChat שלך לא נמצאים לא בחשבון המותג ולא בחשבון Google, רק בשרתים של WeChat עצמה.
  \item \textbf{לכן כשמשתמש אנדרואיד עושה גיבוי, צריך לבדוק בנפרד: גיבוי ענן מותג, גיבוי Google, גיבוי היסטוריית צ'אט WeChat - שלושתם חיוניים.}שש, שיתוף משפחתי: גם בני המשפחה נהנים
\end{itemize}

“שיתוף משפחתי” של Apple מאפשר לך לשתף עם המשפחה רכישות אפליקציות, מוזיקה, סרטים, מקום iCloud, וגם מיקום. עד 6 אנשים. שיטת ההגדרה: הגדרות השם שלך שיתוף משפחתי. ל-Huawei ול-Xiaomi יש פונקציות דומות, ניתן למצוא “שיתוף משפחתי” או “ענן משפחתי” בהגדרות החשבון.

\booksection{תרחיש שימושי מאוד: הגדר שיתוף משפחתי בטלפון של הוריך, כך תוכל לראות את המיקום שלהם בכל עת (בתנאי שהם מסכימים), למקרה שקשיש ילך לאיבוד או לא ניתן יהיה להשיג אותו בטלפון, יהיה מיקום.}

השלכות פריצה לחשבון חמורות ממה שתדמיין

אם האקר ישיג את סיסמת החשבון שלך, הוא יכול: לנעול את הטלפון מרחוק ולדרוש כופר, לצפות בתמונות ובאנשי הקשר שלך, לקנות בשיטת התשלום שלך. יותר גרוע, אם החשבון שלך קשור ליותר מדי דברים, שחזור החשבון עצמו הוא סיוט. לכן, אימות דו-שלבי - לך להפעיל עכשיו.

\begin{tcbwarn}{תרגול מעשי: בדיקת רפואה מקיפה לחשבון שלך}
פתח את הגדרות הטלפון, כנס לדף החשבון, אשר שאימות דו-שלבי מופעל.

\end{tcbwarn}

\begin{tcbpractice}{בדוק ש"מספר טלפון מהימן" ו"מפתח שחזור" הוגדרו, המפתח רשום על נייר.}
\begin{enumerate}
  \item כנס לניהול iCloud או שטח ענן, ראה אילו נתונים מסונכרנים, כבה את מה שלא צריך.
  \item אשר ש"איתור הטלפון שלי" מופעל.
  \item אם אתה משתמש באנדרואיד, בדוק שגיבוי ענן מותג, גיבוי Google, גיבוי WeChat - כולם מכוסים.
  \item פרק 17: קישוריות הטלפון עם מכשירים אחרים
  \item הטלפון בידך הוא למעשה שלט רחוק על ומרכז נתונים. הוא לא רק מתקשר, אלא גם יכול לקשור יחד מחשב, טלוויזיה, שעון, רמקול, מכונית, ולהפוך את המכשירים האלה לאיים לא מבודדים. בפרק זה נדבר איך לגרום להם "לשחק כצוות".
\end{enumerate}

\end{tcbpractice}
